% ---------------------------------------------------
% Trabajo Final de Grado
% Author: Fabián González Lence <alu0101549491@ull.edu.es>
% Chapter: Related Work
% ----------------------------------------------------

\cleardoublepage
\chapter{Estado del Arte} \label{chap:RelatedWork}

La integración de la \ac{IA} en la ingeniería del software ha evolucionado rápidamente con la aparición de los \ac{LLMs}. Diversas investigaciones recientes \cite{Wang:2023:LLMReview, Tosi:2024:CodeQuality, Turunen:2023:FromIdeas, Chen:2021:Codex, Khade:2025:AIinSD, Qian:2024:ChatDev} coinciden en que estos modelos están transformando la forma en que se concibe, desarrolla y mantiene el software. El presente capítulo revisa los trabajos más relevantes en torno a tres ejes temáticos: la generación de código mediante \ac{LLMs}, la evaluación de la calidad del software producido por \ac{IA}, y los enfoques emergentes de colaboración multiagente.

\section{Generación de código con LLMs} \label{sec:generacion-codigo}

La capacidad de los \ac{LLMs} para generar código a partir de descripciones en lenguaje natural constituye una de las aplicaciones más estudiadas de la \ac{IA} generativa en ingeniería del software. Modelos como GPT-4 (OpenAI)\footnote{GPT-4: \url{https://openai.com/index/gpt-4/}}, Claude (Anthropic)\footnote{Claude: \url{https://claude.com/product/overview}}, Codex (OpenAI)\footnote{Codex: \url{https://openai.com/codex/}}, CodeT5 (Salesforce)\footnote{CodeT5: \url{https://github.com/salesforce/codet5}} o StarCoder (BigCode)\footnote{StarCoder: \url{https://huggingface.co/blog/starcoder}} se han integrado progresivamente en los entornos de desarrollo, ofreciendo asistencia en tiempo real para la escritura, completado y refactorización de código.\\

Wang y Chen \cite{Wang:2023:LLMReview} presentan una revisión del panorama actual de la generación de código con \ac{LLMs}, destacando su notable capacidad para producir código funcional en múltiples lenguajes. Su análisis abarca las principales arquitecturas y \textit{benchmarks} empleados, aunque subrayan que la evaluación de la calidad del código generado no ha avanzado al mismo ritmo que las aplicaciones prácticas.\\

El \textit{benchmark} más influyente en este dominio es \textbf{HumanEval} \cite{Chen:2021:Codex}, introducido junto con Codex ---el modelo subyacente a GitHub Copilot--- y diseñado para medir la corrección funcional mediante la métrica \textit{pass@k}. Sus limitaciones ---como problemas autocontenidos, dominio reducido o ausencia de contexto arquitectónico--- son precisamente las que motivan el enfoque de este trabajo, centrado en proyectos completos en lugar de fragmentos aislados.\\

A pesar del impacto innegable de estas herramientas, la generación de código mediante \ac{LLMs} presenta limitaciones bien documentadas. Los modelos tienden a producir código que funciona para los casos de prueba más evidentes pero que puede fallar en casos límite, presentar vulnerabilidades de seguridad o no adherirse a los principios de diseño del proyecto. Además, la calidad del resultado depende en gran medida de la calidad del \textit{prompt}: especificaciones ambiguas o incompletas conducen invariablemente a código que requiere intervención humana significativa.

\section{Evaluación de la calidad del código generado por IA} \label{sec:evaluacion-calidad}

Si bien la capacidad de los \ac{LLMs} para producir código funcional ha quedado ampliamente demostrada, la \textbf{calidad} de ese código más allá de su corrección inmediata es una cuestión considerablemente menos resuelta.\\

Tosi \cite{Tosi:2024:CodeQuality} realiza una de las evaluaciones empíricas más completas hasta la fecha, comparando la calidad del código generado por GPT-3.5, GPT-4 y Google Bard en un conjunto diverso de tareas. Sus resultados confirman que, aunque estos modelos producen código funcional en la mayoría de los casos, presentan limitaciones en mantenibilidad y adherencia a buenas prácticas: el código tiende a exhibir \textbf{complejidad ciclomática moderada-alta} y \textit{code smells} recurrentes ---métodos excesivamente largos, duplicación de lógica o acoplamiento innecesario---. Tosi concluye que la correcta definición de requisitos en el \textit{prompt} es determinante para la calidad del resultado, y que la supervisión experta sigue siendo imprescindible.\\

Cabe señalar que la métrica \textit{pass@k} \cite{Chen:2021:Codex} resulta insuficiente como indicador de calidad global: un programa puede superar todos sus casos de prueba y presentar simultáneamente alta complejidad ciclomática o ausencia de principios como \ac{SOLID}, lo que refuerza la necesidad de un análisis multidimensional como el adoptado en este trabajo.\\

Khade y Sambhe \cite{Khade:2025:AIinSD} amplían la perspectiva documentando cómo se pueden emplear herramientas de \ac{IA} para la detección automatizada de defectos, el análisis estático de seguridad y la generación de conjuntos de pruebas. Esta idea ---que la \ac{IA} pueda generar y revisar código--- es central en la metodología de este trabajo, donde se asignan modelos específicos a las tareas de revisión y \textit{testing}, evaluando su capacidad para detectar defectos en código producido por un modelo distinto.

\section{Colaboración multiagente en el desarrollo de software} \label{sec:colaboracion-multiagente}

Más allá del uso de modelos individuales para tareas aisladas, una línea emergente explora la organización de múltiples agentes de \ac{IA} en flujos de trabajo colaborativos que reproduzcan la dinámica de un equipo de desarrollo humano, donde cada agente asume un rol especializado ---analista, arquitecto, desarrollador, revisor, \textit{tester}--- siguiendo las fases del \ac{SDLC}.\\

Turunen y Fagerholm \cite{Turunen:2023:FromIdeas} exploran la automatización en el desarrollo de aplicaciones mediante agentes como AutoGPT\footnote{AutoGPT: \url{https://agpt.co/}} o Microsoft AutoGen\footnote{Microsoft AutoGen: \url{https://www.microsoft.com/research/project/autogen/}}, describiendo escenarios donde distintos agentes generan código, lo revisan, producen pruebas y coordinan el flujo global. Su conclusión principal es que el futuro de la ingeniería del software podría pasar por entornos de agentes colaborativos, aunque reconocen que la coordinación efectiva entre agentes sigue siendo un problema abierto.\\

Dos trabajos recientes ilustran este paradigma con mayor concreción. ChatDev \cite{Qian:2024:ChatDev} organiza agentes especializados en un flujo de diseño, codificación y \textit{testing} guiado por lenguaje natural, completando proyectos simples en menos de siete minutos. MetaGPT \cite{Hong:2024:MetaGPT} va más lejos al codificar \ac{SOPs} en los prompts de cada agente ---gestor de producto, arquitecto, ingeniero, tester--- exigiendo artefactos intermedios que reducen la acumulación de errores.\\

Este enfoque presenta ventajas claras: permite explotar las \textbf{fortalezas complementarias} de distintos modelos, introduce mecanismos implícitos de control de calidad al separar el rol de generación del de revisión, y facilita la trazabilidad del proceso mediante artefactos intermedios. Sin embargo, también plantea desafíos como la \textbf{pérdida de contexto} entre fases, la \textbf{propagación de errores} desde etapas tempranas, y las fricciones por la disparidad de los modelos.\\

Khade y Sambhe \cite{Khade:2025:AIinSD} enfatizan que la \ac{IA} en el desarrollo de software no debe entenderse como un reemplazo del ingeniero humano, sino como un \textbf{aumento de sus capacidades}. El escenario más productivo es aquel en el que la \ac{IA} automatiza las tareas más repetitivas ---generación de código \textit{boilerplate}\footnote{código \textit{boilerplate}: Código plantilla que se reutiliza en distintas partes de un proyecto sin cambios significativos, para establecer la estructura básica o configuración inicial.}, escritura de pruebas, detección de code smells--- mientras el ser humano aporta el juicio crítico y la visión arquitectónica. Esta visión es precisamente la adoptada en este trabajo.

\section{Posicionamiento del presente trabajo} \label{sec:posicionamiento}

La revisión de la literatura revela un campo en rápida expansión con vacíos significativos que este \ac{TFG} pretende abordar. La Tabla~\ref{tab:comparativa-trabajos} sintetiza las principales diferencias entre los trabajos revisados y la contribución propuesta.\\

\begin{table}[htbp]
\centering
\begin{tabular}{lcccc}
\hline
\textbf{Característica} & \textbf{Wang} & \textbf{Tosi} & \textbf{Turunen} & \textbf{Este TFG} \\
 & \cite{Wang:2023:LLMReview} & \cite{Tosi:2024:CodeQuality} & \cite{Turunen:2023:FromIdeas} & \\
\hline
Múltiples modelos de IA          & No  & Sí  & Sí  & Sí \\
Roles especializados por fase    & No  & No  & Sí  & Sí \\
Proyectos completos (no fragmentos) & No  & No  & No  & Sí \\
Complejidad creciente            & No  & No  & No  & Sí \\
Revisión cruzada entre modelos   & No  & No  & No  & Sí \\
Testing generado por IA          & No  & No  & No  & Sí \\
Documentación de prompts         & No  & Parcial & Parcial & Amplia \\
Métricas de calidad de código    & Teórico & Sí  & No  & Sí \\
\hline
\end{tabular}
\caption{Comparativa entre los trabajos revisados y el presente \ac{TFG}.}
\label{tab:comparativa-trabajos}
\end{table}

La investigación existente se ha centrado predominantemente en tareas aisladas evaluadas con \textit{benchmarks} como HumanEval o MBPP \cite{Wang:2023:LLMReview}, dejando sin explorar el desarrollo de aplicaciones completas con arquitecturas reales y requisitos no funcionales. Este \ac{TFG} aborda ese vacío posicionándose en la intersección entre generación de código, evaluación de calidad y colaboración multiagente, evaluando cinco proyectos web completos de complejidad creciente, documentando de forma rigurosa el proceso completo, desde los \textit{prompts} empleados hasta las decisiones tomadas en cada etapa.\\

\begin{comment}
    Reseñar/Revisar aquí algunos de los trabajos más relevantes (related work) que hacen cosas similares a las que nosotros queremos hacer.\\
    
    Algunos de estos trabajos parecen relevantes.
    En cualquier caso, tener en cuenta la fecha de publicación de cada uno de ellos.
    Los más recientes sería conveniente revisarlos y a partir de ellos hallar otros trabajos relacionados.
    
    \begin{itemize}
        \item Maleki \& Zhao (2024) \cite{Maleki:2024:PCG} hacen una revisión reciente sobre generación procedural de contenido en videojuegos.
         Dado lo reciente de este trabajo, yo empezaría por estudiar este.
    
        \item Este trabajo \cite{Latif:2022:CEP} (de 2022) presenta una revisión crítica de herramientas y algoritmos para generación procedural de terrenos.
    
        \item Este \cite{Hendrikx:2013:PCG} (año 2013) es un Survey general sobre generación procedural de contenido en juegos.
    
        \item Este otro \cite{Raffe:2012:SPT} (2012) es otro survey, en este caso específico sobre generación procedural de terreno con algoritmos evolutivos.
    
        \item Otra revisión más \cite{Smelik:2014:SPM} (2014) sobre generación procedural de mundos virtuales.
    \end{itemize}
\end{comment}

\begin{comment}
La integración de la inteligencia artificial en la ingeniería del software ha evolucionado rápidamente con la aparición de los LLMs. Diversas investigaciones recientes coinciden en que estos modelos están transformando la forma en que se concibe, desarrolla y mantiene el software, al permitir automatizar tareas que tradicionalmente requerían intervención humana directa.\\

Wang y Chen \cite{Wang:2023:LLMReview} revisan el panorama actual de la generación de código con LLMs, destacando que estos modelos han demostrado una notable capacidad para comprender y producir código funcional en múltiples lenguajes de programación. Sin embargo, subrayan una carencia importante: la evaluación sistemática de la calidad del código no ha mostrado un avance al mismo ritmo que las aplicaciones prácticas, dejando abierta la necesidad de métodos de análisis más completos que consideren múltiples características de calidad más allá de la corrección funcional.\\

En la misma línea, Turunen y Fagerholm \cite{Turunen:2023:FromIdeas} exploran la automatización de la creación de aplicaciones a partir de ideas o especificaciones textuales, mostrando cómo herramientas basadas en agentes de IA (como AutoGPT o Microsoft AutoGen) permiten distribuir las tareas del ciclo de vida del software entre distintos roles inteligentes que colaboran mediante conversaciones. Estos autores apuntan que el futuro de la ingeniería del software podría pasar por entornos de agentes de IA colaborativos, en los que cada modelo actúe como desarrollador, tester, o gestor de requisitos, transformando radicalmente los flujos de trabajo tradicionales.\\

Por su parte, Davide Tosi \cite{Tosi:2024:CodeQuality} realiza una evaluación empírica de la calidad del código generado por distintos LLMs (GPT-3.5, GPT-4 y Google Bard), concluyendo que, aunque estos modelos logran resultados funcionales y coherentes en la mayoría de los casos de prueba, aún requieren supervisión constante de expertos en software y una correcta definición de requisitos para alcanzar una calidad excepcional. Además, observa que el código generado tiende a presentar una complejidad moderada-alta, junto con algunos code smells, lo que subraya la importancia de la revisión humana, incluso cuando el código funciona correctamente.\\

De modo complementario, Khade y Sambhe \cite{Khade:2025:AIinSD} amplían la perspectiva al examinar la aplicación de la IA no solo en la generación de código, sino también en pruebas automatizadas, mantenimiento y seguridad. Enfatizan que, más que reemplazar a los ingenieros de software, la IA está aumentando sus capacidades al automatizar tareas tediosas, permitiéndoles enfocarse en el diseño de alto nivel y la resolución de problemas complejos, aunque esto requiere una cooperación equilibrada entre humanos y máquinas.\\

En conjunto, estos estudios reflejan un campo de rápida expansión, pero todavía fragmentado. Las investigaciones actuales se centran mayoritariamente en el rendimiento individual de los modelos en tareas específicas, mientras que el potencial de colaboración entre múltiples IAs especializadas a lo largo del ciclo del desarrollo de software sigue poco explorado. Precisamente en este vacío se enmarca el presente Trabajo de Fin de Grado, que propone analizar cómo distintos agentes de IA pueden cooperar de forma complementaria para desarrollar software de manera integrada, evaluando la coherencia, calidad y eficacia del producto resultante.
\end{comment}